% Updated from curriculum-octubre2026-final.pdf; original cv_5pages.tex preserved.
% Legacy-only REACT and ENIPROGETTI entries retained from the original CV.
% ANECA project end dates of July 2026 may be reporting cut-offs; no ongoing
% participation is inferred. Employment dates below retain the original CV
% where the ANECA document only gives project/fellowship dates.
\documentclass[a4paper,11pt]{article}
\usepackage[utf8]{inputenc}
\usepackage[T1]{fontenc}
\usepackage[english]{babel}
\usepackage[a4paper,left=2.2cm,right=2.2cm,top=1.7cm,bottom=1.8cm]{geometry}
\usepackage{graphicx,enumitem,xcolor,hyperref,fancyhdr,lastpage}
\usepackage{xurl}
\setlength{\emergencystretch}{2em}
\hypersetup{hidelinks,pdftitle={Laura Moreno Martínez — Curriculum vitae, October 2026},pdfauthor={Laura Moreno Martínez}}
\definecolor{heading}{RGB}{30,55,75}
\pagestyle{fancy}
\fancyhf{}
\fancyfoot[L]{\small Laura Moreno Martínez\enspace |\enspace October 2026}
\fancyfoot[R]{\small\thepage\ /\ \pageref{LastPage}}
\renewcommand{\headrulewidth}{0pt}
\setlength{\parindent}{0pt}
\setlength{\parskip}{5pt}
\setlist[itemize]{leftmargin=*,itemsep=4pt,topsep=3pt,parsep=0pt}
\setlist[enumerate]{leftmargin=*,itemsep=7pt,topsep=4pt,parsep=0pt}
\newcommand{\cvsection}[1]{\par\vspace{8pt}{\large\bfseries\color{heading}#1}\par\vspace{3pt}}
\newcommand{\entry}[3]{\textbf{#1}\hfill\textbf{#2}\\#3\par}
\newcommand{\me}{\textbf{Moreno, L.}}
\begin{document}
\begin{minipage}[c]{0.78\textwidth}
{\LARGE\bfseries Laura Moreno Martínez}\\[5pt]
{\large Curriculum vitae}\quad October 2026
\end{minipage}\hfill
\begin{minipage}[c]{0.17\textwidth}
\includegraphics[width=\linewidth,height=2.7cm,keepaspectratio]{../../../laura_photo2024.jpg}
\end{minipage}

\cvsection{Personal information}
\href{mailto:laura.morenomartinez@ua.es}{laura.morenomartinez@ua.es}\\
ORCID: \href{https://orcid.org/0000-0002-8163-0877}{0000-0002-8163-0877}\quad
Scopus: 57209293676\quad ResearcherID: AFX-6764-2022\\
Date of birth: 16 October 1990

\cvsection{Scientific career description}
Assistant Professor at the University of Alicante, working in \textbf{applied mathematics and computational mechanics}. Her research focuses on stabilized finite element formulations, the Material Point Method (MPM), viscoelastic and viscoplastic flows, fluid--structure interaction, and multiphysics simulation of natural hazards and large deformations.

She obtained her PhD at Universitat Politècnica de Catalunya (UPC) in 2021, supported by a competitive Spanish FPI fellowship, and subsequently held postdoctoral positions at the University of Padua and TU Braunschweig. Her work combines mathematical formulation, numerical analysis, implementation and validation, including contributions to the open-source \textbf{Kratos Multiphysics} software.

Her publication record comprises \textbf{10 journal articles}, including nine in JCR-indexed journals (eight in Q1) and five as first author. She has been awarded one \textbf{six-year research evaluation (sexenio), 2019--2024}. She has experience in competitive national and international projects and industrial research contracts, alongside international teaching and co-supervision of two Master's theses and one Bachelor's thesis.

\cvsection{Current position}
\entry{Assistant Professor}{September 2024--present}{Department of Applied Mathematics and Aerospace Engineering, University of Alicante, Spain.}

\cvsection{Previous positions}
\entry{Postdoctoral researcher}{September 2023--August 2024}{Institute of Structural Analysis, TU Braunschweig, Germany. Particle methods for the constitutive evolution and aging of building materials and structures.}
\entry{Postdoctoral researcher}{October 2021--September 2023}{Department of Mathematics ``Tullio Levi-Civita'', University of Padua, Italy. Finite element and particle methods for coupled multiphysics problems and natural hazards.}
\entry{Predoctoral researcher}{September 2016--September 2021}{CIMNE / UPC, Barcelona, Spain. Stabilized finite element methods for highly elastic viscoelastic flows, including thermal effects; implementation in an object-oriented Fortran HPC code.}
\entry{Solutions Assistant (SA-N1)}{December 2015--August 2016}{Everis (NTT Data), Barcelona, Spain. Technical and functional support for a large database system.}
\entry{Research Assistant}{September 2014--February 2015}{University of Santiago de Compostela, Spain. Numerical simulation of electromagnetic processes with hysteresis.}

\newpage
\begingroup\small
\cvsection{Education}
\entry{PhD in Structural Analysis}{22 September 2021}{Universitat Politècnica de Catalunya, Spain. Excellent \textit{cum laude}.\\
Thesis: \textit{Numerical modelling of viscoelastic flows based on a log-conformation formulation}. Supervisors: Ramon Codina and Joan Baiges.}
\entry{Master's degree in Industrial Mathematics}{2013--2015}{University of Santiago de Compostela, Spain. Specialization in numerical methods and finite element simulation of industrial processes.}
\entry{Degree in Mathematics (five-year Licenciatura)}{2008--2013}{University of Alicante, Spain.}

\cvsection{Fellowships and awards}
\begin{itemize}
\item \textbf{Research sexenio, 2019--2024}, awarded by ANECA.
\item \textbf{DFG-funded postdoctoral fellowship}, TU Braunschweig, 2023--2024.
\item \textbf{Postdoctoral fellowships}, University of Padua, 2021--2023.
\item \textbf{Spanish FPI predoctoral fellowship}, 2017--2021, linked to ELASTIC-FLOW (DPI2015-67857-R).
\item \textbf{PhD awarded Excellent \textit{cum laude}}, 2021.
\end{itemize}

\cvsection{Participation in research projects}
\textbf{National and international projects}
\begin{itemize}
\item \textbf{SEA4FUTURE} --- Sustainable and Ecosystem-based Approach for Fisheries through the Utilization of Technology to Uphold Resilient Environments. European Commission; UA PI: César Bordehore Fontanet. Research collaborator in scientific, training and dissemination activities, documented January--July 2026. Project scheduled to end in December 2029.
\item \textbf{Satellite Geodesy for Earth Observation and Climate Studies (SG), VIGROB-260}. University of Alicante; PI: María Isabel Vigo Aguiar. Research collaborator in training and scientific dissemination activities, January--July 2026.
\item \textbf{GRK 2075} --- Modelling the constitutive evolution of building materials and structures with respect to aging. DFG / TU Braunschweig; PIs: Manfred Krafczyk and Roland Wüchner. 2023--2024. Model development, implementation in Kratos Multiphysics and numerical validation.
\item \textbf{REACT} --- Digital Twins of Civil Structures and Protection Systems in a Climate Change Perspective. TUM Institute for Advanced Study; PI: Antonia Larese. Participation: October 2022--September 2023.
\item \textbf{SID 2020, University of Padua} --- Numerical methods for multiphysics problems. PI: Antonia Larese. 2021--2023. Development and implementation of models for coupled problems and extreme natural events in Kratos Multiphysics.
\item \textbf{ELASTIC-FLOW (DPI2015-67857-R)} --- Efficiency optimization in mixing and heat transfer using viscoelastic fluids. MINECO; PIs: Ramon Codina and Joan Baiges. FPI-linked participation, 2017--2021. Stabilized finite element formulations, implementation and validation.
\item \textbf{FORJACEMIC}. CIE-GALFOR (Innterconecta); PI: Alfredo Bermúdez. September--December 2014. Numerical simulation of electromagnetic processes in automotive crankshaft manufacturing.
\end{itemize}
\textbf{Industrial research and consulting contracts}
\begin{itemize}
\item \textbf{ENIPROGETTI / ENI}, 2022. Sensitivity analysis of a physical-channel model to material parameters and well geometry. PI: Mario Putti.
\item \textbf{Instituto de Astrofísica de Canarias / CIMNE}, March--August 2020. Numerical assessment of the future European Solar Telescope's impact on the optical quality of the William Herschel Telescope. PI: Ramon Codina.
\item \textbf{EWE+ / CIMNE}, March--August 2020. Numerical assessment of aerodynamic loads on the Mont-Royal tunnel separation wall. PIs: Ramon Codina and Joan Baiges.
\end{itemize}

\endgroup
\newpage
\cvsection{Publications}
\textit{Complete journal publication list, in reverse chronological order.}
\begin{enumerate}
\item Rodríguez-Velázquez, M. Á., \textbf{Moreno-Martínez, L.}, Chávez-Negrete, C., Domínguez-Sánchez, C., and Domínguez-Mota, F. J. (2026). Rainfall-Triggered debris flows: The July 10, 2018, event at Cerro del Quinceo and its simulation using the material point method. \textit{Ingeniería Investigación y Tecnología}, XXVII, 1--9.
\item Castañar, I., \me, and Codina, R. (2026). Study of Stabilized Mixed Formulations for Fluid-Structure Interaction Problems within a Variational Multiscale Framework. \textit{Computer Methods in Applied Mechanics and Engineering}, 541, 118648. \href{https://doi.org/10.1016/j.cma.2025.118648}{\nolinkurl{doi:10.1016/j.cma.2025.118648}}.
\item \me, Larese, A., and Wüchner, R. (2025). A mixed stabilized MPM formulation for incompressible hyperelastic materials using Variational Subgrid-Scales. \textit{Computer Methods in Applied Mechanics and Engineering}, 435, 117621. \href{https://doi.org/10.1016/j.cma.2024.117621}{\nolinkurl{doi:10.1016/j.cma.2024.117621}}.
\item Codina, R., Baiges, J., Castañar, I., Martínez-Suárez, I., \me, and Parada, S. (2023). An embedded strategy for large scale incompressible flow simulations in moving domains. \textit{Journal of Computational Physics}, 488, 112181. \href{https://doi.org/10.1016/j.jcp.2023.112181}{\nolinkurl{doi:10.1016/j.jcp.2023.112181}}.
\item \me, Castañar, I., Baiges, J., Codina, R., and Cattoni, D. (2023). Numerical simulation of Fluid-Structure Interaction problems with viscoelastic fluids using a log-conformation reformulation. \textit{Computer Methods in Applied Mechanics and Engineering}, 410, 115986. \href{https://doi.org/10.1016/j.cma.2023.115986}{\nolinkurl{doi:10.1016/j.cma.2023.115986}}.
\item \me, Codina, R., and Baiges, J. (2021). Numerical simulation of non-isothermal viscoelastic fluid flows using a VMS stabilized finite element formulation. \textit{Journal of Non-Newtonian Fluid Mechanics}, 296, 104640. \href{https://doi.org/10.1016/j.jnnfm.2021.104640}{\nolinkurl{doi:10.1016/j.jnnfm.2021.104640}}.
\item Castillo, E., \me, Baiges, J., and Codina, R. (2021). Stabilised Variational Multi-scale Finite Element Formulations for Viscoelastic Fluids. \textit{Archives of Computational Methods in Engineering}, 28, 1987--2019. \href{https://doi.org/10.1007/s11831-020-09526-x}{\nolinkurl{doi:10.1007/s11831-020-09526-x}}.
\item Codina, R., and \me\ (2021). Analysis of a stabilized finite element approximation for a linearized logarithmic reformulation of the viscoelastic flow problem. \textit{ESAIM: Mathematical Modelling and Numerical Analysis}, 55, S279--S300. \href{https://doi.org/10.1051/m2an/2020038}{\nolinkurl{doi:10.1051/m2an/2020038}}.
\item \me, Codina, R., and Baiges, J. (2020). Solution of transient viscoelastic flow problems approximated by a term-by-term VMS stabilized finite element formulation using time-dependent subgrid-scales. \textit{Computer Methods in Applied Mechanics and Engineering}, 367, 113074. \href{https://doi.org/10.1016/j.cma.2020.113074}{\nolinkurl{doi:10.1016/j.cma.2020.113074}}.
\item \me, Codina, R., Baiges, J., and Castillo, E. (2019). Logarithmic conformation reformulation in viscoelastic flow problems approximated by a VMS-type stabilized finite element formulation. \textit{Computer Methods in Applied Mechanics and Engineering}, 354, 706--731. \href{https://doi.org/10.1016/j.cma.2019.06.001}{\nolinkurl{doi:10.1016/j.cma.2019.06.001}}.
\end{enumerate}

\newpage
\cvsection{Conferences and invited talks (selected, 2023--2026)}
\begin{itemize}
\item \textbf{HYDROROM Workshop, Padua, February 2026.} \textit{Multiphysics Modeling of Natural Hazards: Numerical Approaches for Structural Impact Assessment.}
\item \textbf{RSME Biennial Congress, Alicante, January 2026.} \textit{Simulación numérica de flujos viscoplásticos mediante formulaciones mixtas con MPM en contextos hidrológicos extremos.} With A. Larese and R. Wüchner.
\item \textbf{PARTICLES 2025, Barcelona, October 2025.} \textit{A New Stabilized Mixed MPM Framework for Incompressible Viscoplastic Fluids.} With A. Larese and R. Wüchner.
\item \textbf{ECCOMAS 2024, Lisbon, June 2024.} \textit{Stabilized mixed formulation for an implicit Material Point Method for viscoplastic fluids.} With A. Larese, A. Contri and R. Wüchner.
\item \textbf{COEFIS XVI, La Laguna, April 2024.} Invited plenary: \textit{How to simulate extreme water hazard events.}
\item \textbf{PARTICLES 2023, Milan, October 2023.} \textit{Stabilized mixed formulation for incompressible materials by using VMS in a Material Point Method framework.} With A. Larese and R. Wüchner.
\item \textbf{ICIAM 2023, Tokyo, August 2023.} \textit{Numerical modeling of viscoelastic flows with high elasticity.} With J. Baiges and R. Codina.
\end{itemize}

\cvsection{Scientific service and software}
Co-organizer of minisymposium MS08, \textit{Particle-based methods in applied and industrial sciences}, SIMAI 2023, Matera, Italy. Peer reviewer for computational mechanics journals. Contributor to Kratos Multiphysics through model and element implementation, validation and testing.

\cvsection{Memberships of scientific societies}
INDAM--GNCS (since 2023); Spanish Association for Numerical Methods in Engineering, SEMNI (since 2019); National Association of Mathematics Students, ANEM (since 2011).

\newpage
\cvsection{Teaching}
More than \textbf{500 hours across nine undergraduate and Master's courses}, taught in Spanish, English and Italian. Student evaluations at the University of Alicante range from 8.61 to 9.73 out of 10.
\begin{itemize}
\item \textbf{University of Alicante, 2024--2026:} \textit{Fundamentos Matemáticos en Ingeniería I} (Aerospace Engineering; 253 h); \textit{Matemáticas Básicas} (Telecommunication Sound and Image Engineering; 110 h); \textit{Fundamentos de Matemática Aplicada 2} (Robotics Engineering; 16 h); \textit{Fundamentos Matemáticos 2} (Architecture; 22 h). Theory and practical teaching in Spanish; accumulated teaching hours.
\item \textbf{Universitat Oberta de Catalunya, 2026:} \textit{Métodos Numéricos en Ciencia de Datos}, online teaching in the Bachelor's degree in Data Science (60 h).
\item \textbf{Technical University of Munich, 2026:} Guest lecture, \textit{Particle Based Methods in Engineering}, M.Sc. Civil Engineering (2 h, English).
\item \textbf{TU Braunschweig, 2024:} \textit{Particle Methods}, Master's in Civil Engineering (10 h, English).
\item \textbf{University of Padua, 2023:} \textit{Calcolo Numerico}, Chemical and Materials Engineering (24 h, Italian); \textit{Numerical Methods for Continuous Systems}, Master's in Mathematical Engineering (12 h, English).
\end{itemize}
Participation in two UA teaching innovation projects (2025--2026), on self-assessment tools in numerical calculus and student participation and attainment in engineering. Coordination activities in \textit{Yo te ayudo con las mates} (2024--2025).

\cvsection{Supervision}
\begin{itemize}
\item \textbf{Master's thesis, Giovanni Campanella, Padua, July 2025:} \textit{A mixed Material Point Method for the Stokes equations: a stabilized velocity-pressure formulation}. Co-supervised with Antonia Larese.
\item \textbf{Master's thesis, Alessandro Contri, Padua, December 2021:} \textit{An Updated Lagrangian displacement-based formulation for free surface incompressible fluids using MPM}. Co-supervised with Antonia Larese.
\item \textbf{Bachelor's thesis, Lucas Peydró Ferrando, Alicante, July 2026:} \textit{Simulación numérica de las olas generadas por deslizamientos mediante el método del punto material: aplicación al caso de Lituya Bay}. Co-supervisor, Degree in Physics.
\item \textbf{Research internship, Arthur D'Hebrail, Alicante, June--August 2026:} Technical guidance and supervision of training activities.
\end{itemize}

\cvsection{Language skills}
Spanish: native. English, Italian and Catalan/Valencian: used in academic or professional settings; certified B1 in each. English: additional B2.2-level course training. German: certified A1.
\end{document}
